\subsection{Założenia teorytyczne}
\subsubsection{Translacja reprezentacji latentnych modeli do wspólnej przestrzeni}
Istniejące techniki linear translation, dodatkowe bardziej zaawansowane
\subsubsection{Modele multimodalne i wymuszenie semantycznych reprezenacji danych}
Contrastive, triplet loss
\subsection{Metodologia}
\subsubsection{Wykorzystane modele ASR i embeddingowe}
Powtorzenie wnioskow z poprzedniego rozdzialu, dobór modelu embeddingowego
\subsubsection{Model agregujący atencję \textit{Attention Pool Model} wykorzystujacy uczenie nadzorowane}
Omówienie architektury, motywacja, przebieg treningu
\subsubsection{Model oparty o translację reprezentacji latentnych \textit{Semantic Alignment Model} wykorzystujący uczenie nienadzorowane}
Omówienie architektury, motywacja, przebieg treningu
\subsubsection{Wykorzystane techinki uczenia nienadzorowanego Data2Vec we wstępnym treningu}
Omówienie wykorzystania Data2Vec dla każdego z powyższych modeli
\subsubsection{Ewaluacja proponowanych modeli}
Metryki i sposoby ewaluacji pozwalające na jednorodne porównanie modeli.
\subsection{Eksperymenty}
\subsubsection{Założenia projektowe}
Wykorzystane zbiory danych, warianty modeli
\subsubsection{Implementacja}
Opis wykorzystanych narzędzi
\subsubsection{Analiza wyników}
Porównanie stabilności treningu, wpływ Data2Vec na jakość modeli
\subsection{Wnioski}
\subsubsection{Ewaluacja architektury \textit{Attention Pool Model}}
Omówienie, porównanie wariantów, wnioski dotyczące zachowania i generalizacji
\subsubsection{Ewaluacja architektury \textit{Semantic Alignment Model}}
Omówienie, porównanie wariantów, wnioski dotyczące zachowania i generalizacji
\subsubsection{Potencjalne osiągi modeli w kontekście pełnego systemu wyszukiwania}
Zarysowanie powyższych wniosków w realiach systemu wyszukiwania